\subsection{Реализация: Алгоритм Лемке-Хаусона}

Рассмотрим алгоритм Лемке-Хаусона – один из наиболее фундаментальных комбинаторных методов поиска равновесия Нэша в смешанных стратегиях. Ключевой особенностью данного алгоритма является его специализация на биматричных играх, что ограничивает область его применения системами, состоящими ровно из двух агентов. Несмотря на это ограничение, алгоритм Лемке-Хаусона занимает центральное место в вычислительной теории игр, так как он математически гарантирует нахождение как минимум одной точки равновесия в любой конечной игре, предоставляя более эффективную альтернативу полному перебору носителей. Принцип работы метода основан на геометрическом подходе и технике последовательного замещения, родственной симплекс-методу линейного программирования. Алгоритм осуществляет целенаправленный обход вершин многогранников наилучших ответов (best-response polytopes), используя систему числовой маркировки (labeling). Движение по ребрам этих многомерных фигур продолжается до тех пор, пока не будет найдена полностью маркированная пара вершин, координаты которой после нормализации соответствуют искомому профилю равновесных стратегий.

\subsubsection{Математическое обоснование}

Рассмотрим биматричную игру двух лиц. Пусть первый игрок имеет множество чистых стратегий $S_1$ мощности $m$, а второй игрок – множество стратегий $S_2$ мощности $n$. Их выигрыши задаются платежными матрицами $A = (a_{ij})$ и $B = (b_{ij})$ размерности $m \times n$. Смешанные стратегии игроков, представляющие собой вероятностные распределения, обозначим как векторы $x^0 = (x_1^0, \dots, x_m^0)$ и $y^0 = (y_1^0, \dots, y_n^0)$.

Для корректной работы алгоритма Лемке-Хаусона требуется, чтобы все элементы платежных матриц были строго положительными. Для этого вводится константа сдвига $\alpha > \max_{i,j}(a_{ij}, b_{ij})$, которая позволяет перейти к модифицированным матрицам с элементами \cite{nabatova2015}:
\begin{equation}
a'_{ij} = \alpha - a_{ij} > 0, \quad 1 \leqslant i \leqslant m, \quad 1 \leqslant j \leqslant n,
\end{equation}
\begin{equation}
b'_{ij} = \alpha - b_{ij} > 0, \quad 1 \leqslant i \leqslant m, \quad 1 \leqslant j \leqslant n.
\end{equation}

Линейный сдвиг выигрышей не изменяет стратегическую структуру игры и множество её равновесий, однако алгоритмически меняет парадигму с максимизации выигрыша на минимизацию упущенной выгоды. Задача поиска равновесия по Нэшу $(x^0, y^0)$ сводится к геометрическому поиску пары векторов немасштабированных вероятностей $x^* = (x_1^*, \dots, x_m^*)$ и $y^* = (y_1^*, \dots, y_n^*)$, которые формируют многогранники наилучших ответов (best-response polytopes) и удовлетворяют следующим ограничениям:
\begin{equation}
x_i^* \geqslant 0, \quad 1 \leqslant i \leqslant m,
\end{equation}
\begin{equation}
y_j^* \geqslant 0, \quad 1 \leqslant j \leqslant n,
\end{equation}
\begin{equation}
\sum_{j=1}^n a'_{ij} \cdot y^*_j \geqslant 1, \quad 1 \leqslant i \leqslant m,
\end{equation}
\begin{equation}
\sum_{i=1}^m b'_{ij} \cdot x^*_i \geqslant 1, \quad 1 \leqslant j \leqslant n.
\end{equation}

Неравенства (4) и (5) ограничивают область допустимых значений неотрицательным ортантом, в то время как неравенства (6) и (7) задают гиперплоскости, отсекающие начало координат и формирующие замкнутые многогранники \cite{nabatova2015}. 

Ключевым математическим свойством равновесия в алгоритме является условие дополняющей нежесткости (complementary slackness). Оно означает, что рациональный игрок использует чистую стратегию с вероятностью, строго большей нуля, только в том случае, если она является наилучшим ответом на стратегию оппонента (т.е. когда соответствующее ограничение-неравенство обращается в точное равенство единице). Данное условие жестко связывает переменные и формализуется следующими уравнениями:
\begin{equation}
\sum_{i=1}^m x^*_i \cdot \left( \sum_{j=1}^n a'_{ij} \cdot y^*_j - 1 \right) = 0,
\end{equation}
\begin{equation}
\sum_{j=1}^n y^*_j \cdot \left( \sum_{i=1}^m b'_{ij} \cdot x^*_i - 1 \right) = 0.
\end{equation}

Поскольку все $x^*_i$ и $y^*_j$ неотрицательны, а выражения в скобках на основе условий (6) и (7) также неотрицательны, уравнения (8) и (9) выполняются тогда и только тогда, когда в каждом произведении хотя бы один из множителей равен нулю. Алгоритм Лемке-Хаусона осуществляет целенаправленный поиск такой полностью маркированной пары векторов $(x^*, y^*)$, перемещаясь по ребрам многогранников от одной вершины к другой. Геометрически этот обход реализуется через поочередную смену активных ограничений (ярлыков) в виде следующей последовательности шагов \cite{nabatova2015}:

\begin{algorithm}[H]
\caption{Алгоритм Лемке-Хаусона: базовая схема итераций}
\begin{algorithmic}[1]
    \State \textbf{Инициализация:} Стартуя из полностью маркированного начала координат, алгоритм инициирует движение, преднамеренно <<сбрасывая>> один из выбранных ярлыков.
    \State \textbf{Движение по ребру:} Алгоритм продвигается вдоль образовавшегося одномерного ребра внутри одного из многогранников до пересечения с новой гиперплоскостью.
    \State \textbf{Обновление маркеров:} При достижении этой новой вершины алгоритм <<подбирает>> соответствующий ей новый ярлык.
    \State \textbf{Переключение (обработка дубликатов):} Если подобранный ярлык уже присутствует у текущей вершины в парном многограннике (возникает дубликат), алгоритм переключается на второй многогранник. Он возобновляет движение уже в нем, отходя от гиперплоскости, соответствующей дублирующемуся ярлыку.
    \State \textbf{Условие остановки:} Данный альтернирующий процесс продолжается строго до тех пор, пока на очередном шаге не будет подобран тот самый ярлык, который был сброшен в самом начале. Восстановление полного уникального набора ярлыков гарантирует нахождение искомой точки равновесия.
\end{algorithmic}
\end{algorithm}

После того как немасштабированные векторы $x^*$ и $y^*$, удовлетворяющие всем вышеописанным условиям, найдены, итоговые равновесные стратегии Нэша $(x^0, y^0)$ вычисляются путем стандартной нормализации, что возвращает векторы в симплекс вероятностей:
\begin{equation}
x_i^0 = \frac{x_i^*}{\sum_{i=1}^m x_i^*}, \quad 1 \leqslant i \leqslant m,
\end{equation}
\begin{equation}
y_j^0 = \frac{y_j^*}{\sum_{j=1}^n y_j^*}, \quad 1 \leqslant j \leqslant n.
\end{equation}

\subsubsection{Анализ алгоритма}

\textbf{Вычислительная разрешимость и корректность.} Алгоритм является точным комбинаторным методом и математически гарантирует нахождение как минимум одного смешанного равновесия Нэша для любой конечной биматричной игры. Метод всегда завершает свою работу за конечное число шагов. Однако его существенным ограничением является то, что он находит лишь \textit{одно} равновесие, зависящее от выбора стартового ярлыка, и не предоставляет механизма для поиска абсолютно всех точек равновесия в игре.

\textbf{Асимптотическая сложность и эффективность.} 
\begin{itemize}
    \item Средний и лучший случаи: На практике, для случайно сгенерированных платежных матриц, алгоритм часто демонстрирует высокую эффективность, находя решение за время, близкое к полиномиальному.
    \item Худший случай: Долгое время алгоритм считался полиномиальным, однако в 2004 году Савани и фон Штенгель строго доказали, что существуют классы игр, для которых количество шагов растет экспоненциально в зависимости от размерности игры, достигая асимптотической сложности $O(2^n)$.
\end{itemize}

\textbf{Особенности, ограничения и устойчивость.}
Главным структурным ограничением метода является его строгая привязка к биматричным играм. Алгоритм не масштабируется на игры $N$ лиц. 

Особую проблему для вычислительной устойчивости представляют  вырожденные игры. Вырождение возникает, когда в одной вершине многогранника пересекается больше гиперплоскостей, чем размерность пространства (например, когда несколько стратегий приносят одинаковый выигрыш). В таких точках алгоритм может зациклиться, бесконечно меняя базис без реального продвижения. Данная проблема обычно решается методом лексикографического возмущения – добавлением бесконечно малых величин $\epsilon$ к выплатам для искусственного разрушения симметрии.

\textbf{Используемые структуры данных и их влияние.}
\begin{itemize}
    \item Симплекс-таблицы: Основная структура данных, представляющая систему уравнений с остаточными переменными. Память для хранения таблиц требует $O(m \times n)$ пространства.
    \item Множества ярлыков (Labels): Для отслеживания текущих активных ограничений используются булевы массивы или хэш-множества, обеспечивающие доступ и проверку дубликатов за $O(1)$.
    \item Типы данных и потеря точности: Выбор числовых типов оказывает критическое влияние на корректность работы. Поиск следующей вершины осуществляется через тест минимального отношения. Использование стандартных чисел с плавающей точкой может привести к ошибкам округления машинной арифметики. Из-за микроскопических погрешностей алгоритм может неправильно определить минимальное отношение, выбрать неверный ярлык и покинуть разрешенный многогранник, что приведет к сбою (ошибке выполнения). Для создания надежной реализации рекомендуется использовать структуры точной рациональной арифметики, которые сохраняют числитель и знаменатель, гарантируя абсолютную математическую точность при переходах.
\end{itemize}

\subsubsection{Разбор примера работы алгоритма Лемке-Хаусона}

Для демонстрации работы алгоритма рассмотрим игру размерности $3 \times 2$.
Первый игрок имеет множество чистых стратегий $S_1 = \{1, 2, 3\}$, а второй игрок – $S_2 = \{4, 5\}$. 

Исходные платежные матрицы $A$ и $B$ заданы следующим образом:
\begin{equation}
A = \begin{pmatrix}
    3 & 3 \\ 
    2 & 5 \\
    0 & 6 
    \end{pmatrix}, 
\quad 
B = \begin{pmatrix} 
    1 & 0 \\ 
    0 & 2 \\ 
    4 & 3 
    \end{pmatrix}
\end{equation}

\textbf{Шаг 1: Подготовка матриц и построение неравенств.}
Для корректной работы симплекс-метода внутри алгоритма Лемке-Хаусона необходимо, чтобы все элементы матриц были строго положительными ($> 0$). Поскольку в матрицах присутствуют нули, мы применяем сдвиг, прибавляя константу $\alpha = 1$ ко всем элементам. Стратегическая суть игры при этом не меняется.

\begin{equation}
A' = A + 1 = \begin{pmatrix} 
    4 & 4 \\ 
    3 & 6 \\ 
    1 & 7 \end{pmatrix}, 
    \quad 
B' = B + 1 = \begin{pmatrix} 
    2 & 1 \\ 
    1 & 3 \\ 
    5 & 4 \end{pmatrix}
\end{equation}

Мы используем классическую формулировку ограничений-неравенств вида $\le 1$, что позволяет нам начать поиск из начала координат. Многогранник $P$ для немасштабированных вероятностей первого игрока $(x_1, x_2, x_3)$ ограничен столбцами матрицы $B'$:
\begin{align}
    x_1 \ge 0 \quad (\text{Ярлык } 1) \\
    x_2 \ge 0 \quad (\text{Ярлык } 2) \\
    x_3 \ge 0 \quad (\text{Ярлык } 3) \\
    2x_1 + 1x_2 + 5x_3 \le 1 \quad (\text{Ярлык } 4) \\
    1x_1 + 3x_2 + 4x_3 \le 1 \quad (\text{Ярлык } 5)
\end{align}

Многогранник $Q$ для немасштабированных вероятностей второго игрока $(y_4, y_5)$ ограничен строками матрицы $A'$:
\begin{align}
    y_4 \ge 0 \quad (\text{Ярлык } 4) \\
    y_5 \ge 0 \quad (\text{Ярлык } 5) \\
    4y_4 + 4y_5 \le 1 \quad (\text{Ярлык } 1) \\
    3y_4 + 6y_5 \le 1 \quad (\text{Ярлык } 2) \\
    1y_4 + 7y_5 \le 1 \quad (\text{Ярлык } 3)
\end{align}

Для вычислений переведем неравенства в систему уравнений, добавив остаточные переменные $s_i \ge 0$.

Для $P$: $s_4 = 1 - 2x_1 - 1x_2 - 5x_3$ и $s_5 = 1 - 1x_1 - 3x_2 - 4x_3$.

Для $Q$: $s_1 = 1 - 4y_4 - 4y_5$, $s_2 = 1 - 3y_4 - 6y_5$ и $s_3 = 1 - 1y_4 - 7y_5$.

\textbf{Шаг 2: Итеративный поиск равновесия.}
Алгоритм стартует в искусственном равновесии (в начале координат): $x = (0,0,0)$ с активными ярлыками $\{1, 2, 3\}$ и $y = (0,0)$ с активными ярлыками $\{4, 5\}$. Набор полностью маркирован: $\{1, 2, 3, 4, 5\}$. 

Чтобы найти смешанное равновесие, инициируем алгоритм, сбросив Ярлык 2 (это означает, что первый игрок начинает использовать стратегию 2, переменная $x_2$ растет от нуля).

Итерация 1 (Многогранник $P$): Увеличиваем $x_2$. Проверяем, какая гиперплоскость будет достигнута первой с помощью теста минимального отношения:

Для $s_4 = 0 \Rightarrow x_2 = 1/1 = 1$.

Для $s_5 = 0 \Rightarrow x_2 = 1/3$.

Минимальное отношение равно $1/3$, оно достигается на ограничении $s_5$ (Ярлык 5). Переменная $x_2$ заменяет $s_5$ в базисе.
Новая вершина $P$: $x = (0, 1/3, 0)$. Набор ярлыков $P$: $\{1, 3, 5\}$.
Общий набор: $\{1, 3, 5\} \cup \{4, 5\} = \{1, 3, 4, 5, 5\}$. Появился дубликат – Ярлык 5.

Итерация 2 (Многогранник $Q$): Сбрасываем дубликат 5 во втором многограннике. Это означает рост переменной $y_5$ при $y_4=0$.

Для $s_1 = 0 \Rightarrow y_5 = 1/4 = 0.25$.

Для $s_2 = 0 \Rightarrow y_5 = 1/6 \approx 0.166$.

Для $s_3 = 0 \Rightarrow y_5 = 1/7 \approx 0.142$.

Минимальное отношение $1/7$ у ограничения $s_3$ (Ярлык 3). Переменная $y_5$ заменяет $s_3$.
Новая вершина $Q$: $y = (0, 1/7)$. Набор ярлыков $Q$: $\{3, 4\}$.
Общий набор: $\{1, 3, 5\} \cup \{3, 4\} = \{1, 3, 3, 4, 5\}$. Дубликат – Ярлык 3.

Итерация 3 (Многогранник $P$): Сбрасываем дубликат 3, увеличивая $x_3$. 
Пересчитываем уравнения $P$ с учетом нового базиса ($x_2 = 1/3 - 1/3x_1 - 4/3x_3 - 1/3s_5$):
\begin{equation}
s_4 = 1 - 2x_1 - (1/3 - 1/3x_1 - 4/3x_3 - 1/3s_5) - 5x_3 = 2/3 - 5/3x_1 - 11/3x_3 + 1/3s_5.
\end{equation}
Увеличиваем $x_3$:

Для $x_2 \ge 0 \Rightarrow 1/3 - 4/3x_3 \ge 0 \Rightarrow x_3 \le 1/4 = 0.25$.

Для $s_4 \ge 0 \Rightarrow 2/3 - 11/3x_3 \ge 0 \Rightarrow x_3 \le (2/3)/(11/3) = 2/11 \approx 0.181$.

Минимум достигается на $s_4$ (Ярлык 4). $x_3$ заменяет $s_4$.

Координаты новой вершины: $x_3 = 2/11$. $x_2 = 1/3 - 4/3(2/11) = 1/11$. Вершина $x = (0, 1/11, 2/11)$.
Набор ярлыков $P$: $\{1, 4, 5\}$. Общий набор: $\{1, 4, 5\} \cup \{3, 4\} = \{1, 3, 4, 4, 5\}$. Дубликат – Ярлык 4.

Итерация 4 (Многогранник $Q$): Сбрасываем дубликат 4, увеличивая $y_4$.
Пересчитываем уравнения $Q$ ($y_5 = 1/7 - 1/7y_4 - 1/7s_3$):

$s_1 = 1 - 4y_4 - 4(1/7 - 1/7y_4 - 1/7s_3) = 3/7 - 24/7y_4 + 4/7s_3$.

$s_2 = 1 - 3y_4 - 6(1/7 - 1/7y_4 - 1/7s_3) = 1/7 - 15/7y_4 + 6/7s_3$.

Увеличиваем $y_4$:

Для $s_1 \ge 0 \Rightarrow y_4 \le (3/7)/(24/7) = 1/8 \approx 0.125$.

Для $s_2 \ge 0 \Rightarrow y_4 \le (1/7)/(15/7) = 1/15 \approx 0.066$.

Минимум на $s_2$ (Ярлык 2). $y_4$ заменяет $s_2$.

Координаты: $y_4 = 1/15$. $y_5 = 1/7 - 1/7(1/15) = 2/15$. Вершина $y = (1/15, 2/15)$.
Набор ярлыков $Q$: $\{2, 3\}$. Общий набор: $\{1, 4, 5\} \cup \{2, 3\} = \{1, 2, 3, 4, 5\}$.

Набор снова полностью маркирован! При этом мы подобрали Ярлык 2 – именно тот, который преднамеренно сбросили на самом первом шаге. Алгоритм успешно завершает работу.

\textbf{Шаг 3: Нормализация.}
Найденные немасштабированные векторы: $x^* = (0, \frac{1}{11}, \frac{2}{11})$ и $y^* = (\frac{1}{15}, \frac{2}{15})$.
Вычисляем суммы компонент: $\sum x^*_i = \frac{3}{11}$ и $\sum y^*_j = \frac{3}{15} = \frac{1}{5}$.
Нормализуем векторы, разделив компоненты на их суммы, и получаем итоговые равновесные смешанные стратегии Нэша:
\begin{equation}
x^0 = \left(0, \frac{1}{3}, \frac{2}{3}\right), \quad y^0 = \left(\frac{1}{3}, \frac{2}{3}\right)
\end{equation}

Проверка ожидаемых выигрышей:
Если Игрок 2 использует стратегию $y^0$, ожидаемые выигрыши Игрока 1 (по исходной матрице $A$) составят:

Стратегия 1: $3(1/3) + 3(2/3) = 3$.

Стратегия 2: $2(1/3) + 5(2/3) = 12/3 = 4$.

Стратегия 3: $0(1/3) + 6(2/3) = 12/3 = 4$.

Стратегии 2 и 3 приносят одинаковый максимальный выигрыш (4), что доказывает рациональность использования их смешивания игроком 1 (условие безразличия выполнено). Равновесие найдено абсолютно точно, что полностью совпадает с геометрической интерпретацией на графиках симплексов.

\subsubsection{Вычислительные эксперименты и выводы}

В ходе вычислительных экспериментов, проведенных на базе разработанной программной реализации (C++), алгоритм Лемке-Хаусона был протестирован на случайно сгенерированных биматричных играх различных размерностей, вплоть до масштабных систем $1000 \times 1000$. Результаты тестирования позволили сделать ряд важных практических выводов, сопоставляющих теорию с реальными вычислениями.

Экспериментально подтверждена высокая эффективность алгоритма в среднем случае. Несмотря на то что теоретический анализ предсказывает экспоненциальное время работы в худшем случае со сложностью $O(2^n)$, на случайных распределениях выигрышей алгоритм не демонстрирует комбинаторного взрыва. Для игр размерности $100 \times 100$ равновесие уверенно находится за считанные миллисекунды, требуя лишь нескольких сотен шагов замещения.